%--------------------------------------------------------------------------------------------------
\chapter{Study Area, Data, and Methodology}
%--------------------------------------------------------------------------------------------------

This chapter introduces the data sources, mathematical models, and spatial context used to evaluate flood hazard and damage modeling in this thesis. The flood risk assessment is structured around three components: flood hazard (depth), vulnerability (damage functions), and exposure (property value). These components are defined both conceptually and mathematically, and are later integrated into a unified framework for risk estimation and validation.

%--------------------------------------------------------------------------------------------------
\section{Hazard and Depth}
%--------------------------------------------------------------------------------------------------

Flood hazard refers to the spatial distribution and intensity of flooding under different return periods. In this study, hazard is represented by flood depth, modeled as a spatial function for given exceedance scenarios.

\subsection{Mathematical Definition of Flood Hazard}

Flood hazard is modeled as a spatially varying flood depth field:

\begin{equation}
	h : \Omega \times \mathcal{T} \rightarrow \mathbb{R}_{\geq 0}, \quad (x, y, T) \mapsto h(x, y, T)
	\label{eq:hazard-map}
\end{equation}

where:
\begin{itemize}
	\item \( \Omega \subset \mathbb{R}^2 \) is the spatial domain of interest,
	\item \( \mathcal{T} = \{ T_1, T_2, \dots, T_K \} \) is the set of return periods considered (e.g., 10, 100, 500 years),
	\item \( h(x, y, T) \) is the flood depth (in meters) at location \( (x, y) \) for return period \( T \).
\end{itemize}

Return periods \( T \) are commonly used to express the frequency of extreme events. The corresponding annual exceedance probability \( p \) is given by:

\begin{equation}
	p = \frac{1}{T}
	\label{eq:return-period}
\end{equation}

Flood hazard models such as those used in this study estimate water depths \( h(x, y, T) \) based on hydrological and hydraulic simulations. In practice, these functions are discretized into raster data aligned to a grid. The continuous function \( h(x, y, T) \) is approximated by an array \( H_{i,j,k} \), where:
\begin{itemize}
	\item \( i, j \) index the spatial grid cells,
	\item \( k \) indexes the return period \( T_k \in \mathcal{T} \).
\end{itemize}

\begin{figure}[htb]
	\centering
	\includegraphics[width=0.75\textwidth]{figures/hazard_placeholder.pdf}
	\caption{Example of a flood hazard depth map \( h(x, y, T) \) for return period \( T = 100 \) years.}
	\label{fig:hazard-map-placeholder}
\end{figure}


\subsection{Hazard Data Sources}

Two flood hazard models are evaluated and compared in this thesis:

\begin{itemize}
	\item \textbf{WRI Aqueduct River Flood Model:} This is a global riverine flood hazard model developed by the World Resources Institute. It offers flood depth estimates at a resolution of approximately 1 km for a broad range of return periods (from 2 to 1000 years). The model integrates hydrological simulations with elevation data and accounts for future climate scenarios through Representative Concentration Pathways (RCP4.5 and RCP8.5). It is distributed in GeoTIFF raster format using EPSG:4326.

	\item \textbf{Integralna karta globin (IKG):} This is a high-resolution local flood hazard dataset prepared by Slovenian authorities for selected municipalities. The model is based on hydraulic simulations of flood dynamics using fine-scale terrain data and site-specific conditions. IKG maps are originally provided as vector files (shapefiles or GeoPackages), categorized into three inundation classes—minor (0.5 m), significant (1.5 m), and severe (2.5 m)—and are available for return periods of 10, 100, and 500 years. The coordinate reference system is EPSG:3794.
\end{itemize}

For the purposes of this study, all hazard layers were harmonized into a common 5 m raster grid, reprojected into EPSG:3794, and stored in Zarr format to enable aligned processing and evaluation. The key differences between the WRI and IKG models are summarized in Table~\ref{tab:hazard_models}.

\begin{table}[htb]
    \centering
    \caption{Comparison of flood hazard models used in the study.}
    \label{tab:hazard_models}
    \begin{tabular}{|p{4.2cm}|p{5.5cm}|p{5.5cm}|}
        \hline
        \textbf{Feature} & \textbf{WRI Aqueduct River Flood} & \textbf{Integralna karta globin (IKG)} \\
        \hline
        Spatial resolution & $\sim$1 km & 5 m \\
        Coverage & Global & 18 regions in Slovenia \\
        Climate scenarios & RCP-aware (e.g., RCP4.5, RCP8.5) & Not scenario-aware \\
        Depth type & Continuous (in meters) & Discrete classes (0.5, 1.5, 2.5 m) \\
        Return periods & 2 to 1000 years & 10, 100, 500 years \\
        File format & Raster (GeoTIFF) & Vector (shapefile, GPKG) \\
        Coordinate system & EPSG:4326 & EPSG:3794 (transformed) \\
        \hline
    \end{tabular}
\end{table}

%--------------------------------------------------------------------------------------------------
\section{Vulnerability and Damage Function}
%--------------------------------------------------------------------------------------------------

Flood vulnerability quantifies the expected damage to an asset given the intensity of the flood at that location. This relationship is typically modeled via a depth--damage function.

\subsection{Mathematical Definition}

A depth--damage function \( D(h) \) maps flood depth \( h \) to the expected fractional loss in asset value:

\begin{equation}
	D : \mathbb{R}_{\geq 0} \rightarrow [0, 1], \quad h \mapsto D(h)
	\label{eq:damage-function}
\end{equation}

The function \( D(h) \) is assumed to be monotonic and piecewise continuous. It starts at a nonzero depth threshold (typically 0.1–0.2 m), increases with water depth, and saturates at a maximum value (e.g., 0.7 to 0.9) depending on asset type.

\begin{figure}[htb]
	\centering
	\includegraphics[width=0.75\textwidth]{figures/vulnerability_placeholder.pdf}
	\caption{Example depth--damage function \( D(h) \) for a residential building.}
	\label{fig:vulnerability-curve-placeholder}
\end{figure}

\subsection{Vulnerability Data Source}

This thesis uses standardized global depth--damage functions published by Huizinga et al.~\cite{huizinga_global_2017}, developed under the auspices of the European Commission Joint Research Centre (JRC). These functions are stratified by building type (e.g., residential, commercial, industrial) and were applied based on available metadata in the Slovenian dataset.

Damage reports from the Administration for Civil Protection and Disaster Relief of Slovenia serve as the primary reference for validation. These records include georeferenced flood locations, reported water depths, and economic losses. When available, building classifications were matched to appropriate JRC vulnerability curves.

%--------------------------------------------------------------------------------------------------
\section{Flood Risk Modeling}
%--------------------------------------------------------------------------------------------------

Flood risk is the expected monetary loss resulting from a flood event. It is computed as the product of hazard, vulnerability, and exposure. This section formalizes how these components are combined.

\subsection{Risk Definition}

The expected loss at a location \( (x, y) \) for return period \( T \) is given by:

\begin{equation}
	R(x, y; T) = D(h(x, y, T)) \cdot E(x, y)
	\label{eq:local-risk}
\end{equation}

where:
\begin{itemize}
	\item \( h(x, y, T) \) is the modeled flood depth (hazard),
	\item \( D(h) \) is the depth--damage function (vulnerability),
	\item \( E(x, y) \) is the value of the asset located at \( (x, y) \) (exposure).
\end{itemize}

To compute the total loss across the study area \( \Omega \), the pointwise losses are integrated:

\begin{equation}
	R_{\text{total}}(T) = \int_{\Omega} D(h(x, y, T)) \cdot E(x, y) \, dx \, dy
	\label{eq:total-risk}
\end{equation}

For rasterized data indexed by \( (i, j) \), this becomes:

\begin{equation}
	R_{\text{total}}(T_k) \approx \sum_{i=1}^{N_x} \sum_{j=1}^{N_y} D(H_{i,j,k}) \cdot E_{i,j}
	\label{eq:total-risk-discrete}
\end{equation}

This formulation aligns with the Physrisk framework, where hazard, vulnerability, and exposure are collocated in gridded format and combined cell-by-cell.

\subsection{Exposure Data}

Exposure refers to the monetary value of buildings and other structures in the floodplain. For the Slovenian study area, exposure estimates were derived from:
\begin{itemize}
	\item Appraised property values from reported flood events,
	\item Official tax assessment records (when available),
	\item Synthetic estimates based on average construction cost per square meter and building type.
\end{itemize}

All exposure data were harmonized onto the same grid as the hazard and vulnerability layers. Where possible, exposure was normalized to match valuation baselines used in the vulnerability function calibration (e.g., building replacement value).

\subsection{Use Cases}

This modeling framework supports the following analyses:

\begin{itemize}
	\item \textbf{Model comparison:} Side-by-side evaluation of WRI vs IKG predictions.
	\item \textbf{Vulnerability validation:} Testing how well JRC depth--damage functions explain observed flood damages in Slovenia.
	\item \textbf{Sensitivity analysis:} Exploring how proximity to water bodies and spatial resolution affect predicted damages.
	\item \textbf{Integration into platforms:} Documenting data onboarding pipelines for use in Physrisk and ESG disclosure workflows.
\end{itemize}

The implementation and experimental evaluation based on this framework are presented in Chapters~4 and 5.
